% Archived 2026-10-16 from Thesis-writeup/Writing/sections/05_setup.tex (old hidden wip block, superseded by the clean rewrite of Section V; Dirk 2026-10-16).
% Old bullets and subsections, hidden from the review build (2026-10-11).
\begin{wip}
\begin{itemize}
\item Generate closed-loop records from the physics baseline plus a truth-only flexible absorber, retaining exact states solely for diagnostics
\item Use multisine excitation that covers the absorber pole and anti-resonance while respecting position, velocity, yaw, and force limits
\item Assign independent random multisine phases to training, validation, and test records so equal spectra do not duplicate trajectories
\item Include multisine-free cycloidal motion profiles matching the measured machine motion to test the operational distribution shift
\item Use only the corrected all-channel Coulomb dataset matched to the frictionless excitation; the earlier mismatched-amplitude merge is superseded
\item The set-valued friction law requires a velocity threshold below which a rail is treated as stuck. This is set to the Stribeck velocity, $v_{s} = 2.25\cdot10^{-3}$ m/s, taken from Lee et al. \cite{lee2020feeddrive}, Table~5, identified on a ball-screw feed drive with a THK SSR20XW guide (Sec.~5.1). Stribeck velocities reported in the literature span $10^{-5}$ to $10^{-2}$ m/s \cite[Sec.~5.2]{lee2020feeddrive}, so the chosen value sits mid-range
\todo{Conflicts with D-204, where the stick band is \texttt{V\_EPS = (cc1+cc2)/m\_total*ts} labelled THEORY: leine1998, about $3.3\cdot10^{-5}$ m/s before the 3 to 10x margin. Confirm which threshold the reported runs use, and whether the second citation (Stribeck range) is the same paper.}
\item State that output noise is currently disabled; add it only if a level and measurement location can be estimated from Jasper, ILC data, or stationary records
\item State the sample rates, record lengths, number of windows, and total independent motion time rather than quoting overlapping-window count as data volume
\item Decision (RK4 value, moved here from Section~II): one reproducibility sentence giving the RK4 step, equal to the sampling period after downsampling the 20~kHz records to 4~kHz, with one substep (entry script: \texttt{fs\_new=4000}, \texttt{up\_sample=1}). Confirm against the \texttt{config.json} of each reported run, since earlier checkpoints may have used another \texttt{up\_sample}
\item Figure (reuse, no separate asset): let the rebuilt training-objective time strip also identify independent phases, short training windows, encoder history, and the long evaluation horizon
\end{itemize}
\todo{Is the hidden linear absorber sufficient simulation complexity, or is a truth-only cubic spring required to demonstrate nonlinear augmentation?}
\todo{discuss noise model and how we filter noise in the setup}

\subsection{Real Telica data and split}
\begin{itemize}
\item Confirm with Jasper and Dragan which reconstructed ETEL setpoint-generator or ILC-reference motions are operationally representative, including minimum/maximum move, long stroke, acceleration, and dwell
\item Use the supervisor-defined split of 11 training, two validation, and two test operating points at 20 kHz
\item Convert rail positions and currents to logical metres and Newtons; form velocity states by central differences because direct velocity is not logged
\item Keep all iterations at one operating point in the same split to prevent identical reference trajectories leaking across sets
\item Describe the $X$ and $Y$ operating-point grid and show whether training covers the interpolation region occupied by both test points
\item Use iteration 0 for controller checks and later ILC iterations for distinct feedforward conditions without treating them as independent trajectories
\item Mark real-data training and held-out evaluation as planned until a final reproducible run and force-calibration convention exist
\item Figure (new): 2D operating-point grid with train, validation, and test locations, operating-range boundary, and trajectory classes; add a 3D or pairwise state-input coverage view only if it reveals coverage hidden by the 2D map
\end{itemize}
\todo{Freeze the approved operational trajectory set with Jasper and Dragan; should direct drive velocity be requested, and does a third coverage coordinate add information beyond the 2D grid?}

\subsection{Training protocol, comparators, and metrics}
\begin{itemize}
\item Match records, neural seeds, windows, optimiser budget, and checkpoint criterion across baseline-only, unconstrained, tangent-OBC, affine-OBC, and black-box arms; treat multisine phase seeds separately
\item Use physical RMS as the primary candidate metric for tracking and a fixed 200--400 ms post-move settling interval; use NRMS, BFR, or spectra only for a distinct supporting question
\item Report physical combination error, projection overlap, added-state ablations, and encoder-init sensitivity alongside predictive metrics
\item Values referenced from Section~III (moved 2026-10-06), each with an honest basis (derived, inherited, tuned, compute-limited): encoder history $n$, window length $n_f$ and stride, number of additional states $n_{\bar x}$, network width and depth, model rate
\item Compared model (moved from Section~III-A, 2026-10-11): for $n_{\bar x}=0$ the neural network only corrects the six baseline states, giving the static parallel augmentation of~\cite[Table~1]{hoekstra2026lfr}; used in the $n_{\bar x}$ comparison (\texttt{THESIS-RESULTS.md} step~4, run R6)
\item Free-run validation (referenced from Section~III-C): whole records, encoder at the record start and $x^c=0$ once, so model error accumulates over the record unlike in a training window; state the record-length-weighted RMS used for checkpoint selection
\item Decision (moved from Section~II-C): joint-estimation arms start from the nominal $\varphi$ of the reference implementation detuned by 10\% per combination (\texttt{combo\_init\_detune}, HEURISTIC, D-191), with the nominal-parameter prior disabled (D-193); state it as a run setting with its reason
\item A parameter prior bounds the drift of the combinations but does not remove the negation flat direction of Section~\ref{sec:obc}; it is off in every arm (D-193)
\item Practical identifiability (referenced from Section~II-C): report the rank of the stacked one-step transition Jacobian with respect to $\varphi$ on the training data (rank 10, nullity 4 in $\theta$ per the reduced-block docstring), from a saved artefact rather than the docstring
\item Scope (from Section~II-C): the structural result concerns the ideal, noise-free baseline model; whether the closed-loop records excite all ten combinations is checked here; the augmented model, where the network can absorb part of the baseline, is treated in Section~\ref{sec:obc}; the encoder does not affect the result, as it only estimates the initial state of a simulation, not the model
\item Value (moved from Section~II-C): the margin $\gamma=0.99$ in \eqref{eq:mdelta_map}, a HEURISTIC margin inside the admissible region (D-229)
\item Metric (moved from Section~II-C): relative combination error in $\theta_{\mathrm{base}}$, with the normalisation of the small signed $m_\Delta$ stated explicitly
\item State per arm whether raw or reduced coordinates are trained; exact-model parameter recovery is an internal implementation check, not a result (moved from Section~II-C)
\item Give the black-box model validation-selected capacity and horizon so it is a capability comparison rather than an artificially weakened baseline
\item Evaluate held-out position/phase interpolation, ASMPT-relevant multisine-free motion, controller transfer, and operating-range extrapolation as separate rungs
\item Predeclare expectation, acceptance criterion, verdict, and mechanism for every result rather than presenting isolated scores
\end{itemize}
\todo{Decide the primary metric, credible controller-transfer arm, and feasible neural-seed count; do not pool initialization, noise, phase, and operating-point uncertainty.}

\subsection{Reproducibility and implementation checks}
\begin{itemize}
\item Verify controller reconstruction, coordinate transforms, normalisation, RK4 transition, and LFR equivalence before training
\item Use the same controller, data conversion, and rollout entry point for training, validation, and checkpoint selection
\item Controller rate (moved from Section~III-C, 2026-10-06): the records are generated with the controller at 20~kHz, the training loop re-discretises it at the 4~kHz model rate (D-141); the residual form of Section~III-C equals the shared-reference loop only for a compatible controller, signals and initial state
\item Compatibility check for that assumption: run the re-discretised controller on the recorded $r-y$ and compare its feedback force with the recorded one; a rollout with $\hat y=y$ gives $\hat u=u$ by construction and checks nothing
\item Controller mismatch: the records use the controller at 20~kHz, training re-discretises it at 4~kHz, so the residual form of Section~\ref{sec:aug_loop} is approximate and part of the measurement noise is not removed
\todo{Missing: size of the controller mismatch. The records use the controller at 20~kHz, training re-discretises it at 4~kHz. Candidate: the compatibility check of the Setup bullet, reported as the force difference against the recorded $u_{\mathrm{act}}$.}
\item Record configuration, software precision, seed, dataset identifier, and checkpoint provenance for every reported arm
\item Keep oracle states and truth parameters out of the trainable pipeline and label every diagnostic that uses them
\item Figure (existing, appendix candidate): use \texttt{controller-chain-v1} only as an audit trail for controller reconstruction and discretisation, not as a main method figure
\end{itemize}
\todo{Which implementation checks belong in the main paper, and which should move to the appendix or repository documentation?}
\end{wip}
